% Arman Pouyaei — Resume
% Single-file LaTeX source. The portfolio site compiles this in the visitor's
% browser (Tectonic/XeTeX via WebAssembly) when they click "Download résumé".
% Section titles: minimal professional style (bold, no rule, no color).
\documentclass[10pt]{article}

\usepackage[utf8]{inputenc} % no-op under XeTeX (UTF-8 is native there)
% T1 font encoding is for pdfLaTeX only. The browser renderer compiles with
% XeTeX, which uses Unicode (TU) fonts natively — requesting T1 there asks
% for EC-encoded metrics the web TeX Live bundle does not ship.
\ifx\XeTeXrevision\undefined
  \usepackage[T1]{fontenc}
\fi
\usepackage[margin=0.5in,top=0.4in,bottom=0.4in]{geometry}
\usepackage{enumitem}
\usepackage{xcolor}
\usepackage{titlesec}
\usepackage{hyperref}

\definecolor{accent}{HTML}{0E6E6B} % deep teal
\hypersetup{colorlinks=true,urlcolor=accent,linkcolor=accent}

% Section titles: minimal professional style — bold, ragged right,
% no rule, no color. Maximum ATS-friendliness.
\titleformat{\section}{
  \vspace{-2pt}\raggedright\large\bfseries
}{}{0em}{}[\vspace{-4pt}]
\titlespacing*{\section}{0pt}{10pt}{6pt}

\setlist[itemize]{leftmargin=1.3em,itemsep=0.1ex,topsep=0.2ex,parsep=0pt}
\setlist[description]{leftmargin=0pt,itemsep=0.25ex,topsep=0.3ex,parsep=0pt,font=\normalfont\bfseries}
\setlength{\parindent}{0pt}
\pagestyle{empty}
\raggedbottom

% \job{title}{org, location}{dates}{subtitle (may be empty)}
\newcommand{\job}[4]{%
  \par\vspace{0.5ex}\noindent\textbf{#1} \hfill #3\\%
  \noindent\textit{#2}%
  \ifx&#4&\else\\[-0.2ex]\noindent{\small\textit{#4}}\fi%
  \vspace{-0.4ex}%
}

% \school{degree}{school}{dates}{note (may be empty)}
\newcommand{\school}[4]{%
  \par\vspace{0.5ex}\noindent\textbf{#1} \hfill #3\\%
  \noindent\textit{#2}%
  \ifx&#4&\else\\[-0.2ex]\noindent{\small #4}\fi%
  \vspace{-0.4ex}%
}

\begin{document}
\fontsize{9}{10.3}\selectfont

% ---------- Header ----------
\begin{center}
  {\large\bfseries ARMAN POUYAEI, PhD}\\[0.5ex]
  {\footnotesize Princeton, NJ \enspace|\enspace (832) 633-4481 \enspace|\enspace
  \href{mailto:arman.pouyaei@gmail.com}{arman.pouyaei@gmail.com} \enspace|\enspace
  \href{https://linkedin.com/in/arman-pouyaei}{linkedin.com/in/arman-pouyaei}}
\end{center}

% ---------- Summary ----------
\section*{Professional Summary}
Senior data engineer and data scientist with 10+ years building production
machine-learning and predictive-modeling systems on large, high-dimensional
data. PhD-trained in statistical inference, Bayesian estimation, and
optimization. Currently leads a global severe-storm hazard modeling program
end to end --- AI-based weather emulation, GPU-accelerated ensemble
forecasting, extreme value theory, physics-based transport and emission
modeling, and forward-looking scenario design --- from solution architecture
through production delivery. Works fluently in Python across HPC and cloud at
petabyte scale, leads and mentors engineering teams, and translates technical
findings into plain-language recommendations through visualization,
prototypes, and executive briefings. 30+ peer-reviewed publications.

% ---------- Skills ----------
\section*{Technical Skills}
\begin{description}
  \item[Statistical \& Predictive Modeling:] Regression, classification,
    clustering, Bayesian inference and state estimation, time-series
    forecasting, extreme value theory (GEV/GPD, peaks-over-threshold) and
    return-period estimation, probabilistic and ensemble forecasting,
    uncertainty quantification, hypothesis testing, constrained optimization.
  \item[Machine Learning \& AI:] PyTorch, JAX, scikit-learn,
    TensorFlow; XGBoost, Random Forest, CNN/UNet, GNN, PINN, neural operators
    and diffusion-based generative models; AI weather emulators; multi-GPU
    distributed training and inference; LLM and NLP tooling, agentic and
    retrieval workflows.
  \item[Physical \& Numerical Modeling:] Lagrangian particle transport
    and dispersion models, physics-based emission schemes, data assimilation
    (3DVAR, EnKF), climate and atmospheric-chemistry models, downscaling, bias
    correction, scenario and climate-pathway analysis.
  \item[Programming \& Data:] Python (NumPy, pandas, Xarray, Dask,
    GeoPandas), R, SQL, Bash, Fortran; Zarr, Parquet, NetCDF, HDF5; relational
    and columnar stores; schema, warehouse, and data-model design.
  \item[Engineering \& Scale:] AWS (Batch, S3, Lambda), Docker,
    SLURM/MPI, CUDA and multi-GPU clusters; production ETL, automated QA/QC,
    anomaly detection, reproducible pipelines.
  \item[Visualization \& Communication:] Tableau, Matplotlib, Plotly,
    Seaborn, interactive HTML dashboards, PowerPoint; prototypes and
    demonstrations built for non-technical stakeholders and executive
    audiences.
\end{description}

% ---------- Experience ----------
\section*{Professional Experience}

\job{Senior Data Engineer (Data Science)}{MSCI Inc., New York, NY}{Aug 2026 -- Present}{Severe-storm hazard modeling, predictive analytics, and AI platform engineering for climate risk. (First Street acquired by MSCI, August 2026.)}
\begin{itemize}
  \item Lead the severe-storm modeling program end to end --- blizzard, dust
    storm, and sandstorm hazards --- owning solution architecture, model and
    tool development, GPU-accelerated training and inference, validation, and
    production delivery.
  \item Built a global ensemble forecasting capability on AI-based weather
    emulators, generating large perturbed-initial-condition ensembles on GPUs
    at a fraction of the cost of physics-based numerical weather prediction,
    then post-processing members into calibrated hazard probabilities and
    spread-based confidence measures.
  \item Applied extreme value theory (GEV and generalized Pareto,
    peaks-over-threshold) to fit the tails of storm severity distributions and
    derive return-period estimates well beyond the observational record, with
    bootstrapped confidence intervals attached to every return level.
  \item Designed a forward-looking scenario framework that projects hazard
    frequency and severity onto future climate pathways, so present-day risk
    and mid- to late-century outcomes can be compared on a consistent,
    auditable basis.
  \item Developed a Lagrangian particle transport model coupled to a
    physics-based dust and sand emission scheme --- friction-velocity
    thresholds with soil, vegetation, and land-surface controls --- to
    simulate emission, long-range transport, and deposition at the spatial
    resolution needed for asset-level exposure.
  \item Build, tune, and validate machine-learning models (PyTorch, JAX,
    scikit-learn) that score risk exposure across hundreds of millions of
    records --- owning feature engineering, training, cross-validation, and
    production deployment end to end.
  \item Set technical direction across the program and review engineering and
    analytical work; keep the team aligned to milestones through regular
    check-ins with engineers and analysts, and report delivery status to
    senior leadership.
  \item Develop AI and LLM-based tooling and agentic workflows that automate
    model diagnostics, output review, and documentation; partner with
    product, research, and client-facing teams to turn open-ended business
    questions into measurable requirements and tangible prototypes.
\end{itemize}

\job{Data Engineer / Research Scientist}{First Street, New York, NY}{May 2025 -- Aug 2026}{Machine learning and data pipeline engineering for global climate risk analytics.}
\begin{itemize}
  \item Designed and deployed end-to-end ML workflows covering data
    preparation, training, validation, and iterative QA/QC for global
    predictive models, including tuning and scaling UNet-based emulators.
  \item Built an automated quality-control and anomaly-detection framework
    that evaluates model outputs against independent benchmarks and surfaces
    failures through dashboards --- cutting review time during global model
    release cycles.
  \item Architected containerized ETL on AWS Batch serving petabyte-scale
    datasets (Zarr, Parquet, NetCDF) to downstream analytics and client
    products.
  \item Converted exploratory notebook research into production-grade,
    reproducible pipelines on SLURM clusters and AWS, working alongside the
    science team to preserve analytical intent.
\end{itemize}

\job{Postdoctoral Research Associate}{NOAA GFDL / Princeton University}{Mar 2023 -- May 2025}{Cooperative Institute for Modeling the Earth System (CIMES), Princeton University.}
\begin{itemize}
  \item Applied statistical and machine-learning analysis to decadal,
    high-dimensional datasets to isolate driver--response relationships hidden
    in noisy signals; findings published in \textit{Geophysical Research
    Letters}.
  \item Led a modeling workstream spanning four research groups and 8+
    scientists --- defining scope, sequencing deliverables, and keeping
    distributed contributors aligned to a shared roadmap.
  \item Built and benchmarked predictive model components against
    independent observational datasets, documented uncertainty for downstream
    users, mentored graduate students, and delivered invited talks at NASA
    GISS and AGU.
\end{itemize}

\job{Postdoctoral Research Fellow}{University of Houston}{Sep 2022 -- Feb 2023}{Bayesian inference, inverse modeling, and predictive analytics.}
\begin{itemize}
  \item Designed a 3DVAR variational inference framework that assimilates
    high-volume observational streams into predictive models --- a
    large-scale optimization and state-estimation problem over sparse,
    heterogeneous data.
  \item Constrained model uncertainty by fusing satellite and in-situ
    measurement sources; authored competitive research proposals to DOE,
    NOAA, NASA, and NSF and supervised undergraduate and graduate
    researchers.
\end{itemize}

\job{Graduate Research Assistant}{University of Houston}{2018 -- 2022}{Air Quality Forecasting and Modeling Lab --- statistical modeling, data fusion, and satellite data analysis.}
\begin{itemize}
  \item Built and open-sourced a Lagrangian particle-trajectory framework
    (C-TRAIL v1.0) that traces how quantities move through large
    spatiotemporal simulations and attributes downstream outcomes to specific
    upstream sources; adopted by external research groups and extended into a
    diagnostic tool that ranks each source region's contribution during
    extreme episodes.
  \item Developed data-fusion operators that reconcile heterogeneous
    observational sources --- satellite, ground-network, and aircraft
    measurements with differing resolution, bias, and coverage --- within a
    Bayesian assimilation framework, producing constrained estimates from
    sparse and inconsistent inputs.
  \item Designed a physics-based scheme for processes below the model's
    resolved grid scale, measurably improving predictive skill (JAMES 2021),
    and co-developed deep-learning surrogates, spatial-imputation models, and
    super-resolution bias correction with the group's machine-learning team.
\end{itemize}

\job{Graduate Research Assistant}{University of Tehran}{2014 -- 2017}{}
\begin{itemize}
  \item Designed optimization and predictive-modeling frameworks for
    renewable energy systems; published in \textit{Energy Conversion and
    Management} and IEEE.
\end{itemize}

% ---------- Education ----------
\section*{Education}

\school{Ph.D., Atmospheric Science}{University of Houston, TX}{2018 -- 2022}{Quantitative focus: Bayesian data assimilation, numerical modeling, statistical analysis of large observational datasets.}
\school{M.Sc., Renewable Energies Engineering}{University of Tehran, Iran}{2013 -- 2016}{}
\school{B.Sc., Mechanical Engineering}{University of Shiraz, Iran}{2008 -- 2012}{}

\par\vspace{0.6ex}
\noindent\textbf{Awards:} Outstanding Graduate Work Scholarship, University of Houston (2021 \& 2022); full scholarships for M.Sc.\ and B.Sc.

% ---------- Leadership ----------
\section*{Leadership, Mentoring \& External Engagement}
\begin{itemize}
  \item \textbf{Team \& Project Leadership:} Coordinate cross-functional
    delivery at MSCI and First Street through regular working sessions where
    engineers and analysts walk through progress and surface blockers
    together. At Princeton/NOAA GFDL, directed a modeling workstream across
    four research groups and 8+ scientists, with collaborators at NASA, NCAR,
    Brookhaven, and PNNL.
  \item \textbf{Mentoring, Teaching \& Funding:} Mentored undergraduate and
    graduate researchers at Princeton/NOAA GFDL and the University of
    Houston, plus two international high school students through the
    Collegiate Mentorship Program (2024). Teaching assistant for Dynamic
    Meteorology and Principles of Atmospheric Science. Authored proposals to
    DOE, NOAA, NASA, and NSF.
  \item \textbf{Communicating to Mixed Audiences:} Invited speaker, NASA GISS
    Seminar Series (2025) and NOAA GFDL Seminar Series (2024). Conference
    presenter at AGU Fall Meetings (2021--2024), AMS Annual Meeting (2024),
    and CMAS (2021).
  \item \textbf{Professional Service \& Outreach:} Lead convener, AGU 2026
    session on aerosol impacts on climate, air quality, and health; convener
    of three AGU sessions across 2024 and 2025. Peer reviewer for
    JGR-Atmospheres, \textit{Atmospheric Environment}, \textit{Scientific
    Reports}, and six additional journals. Organizing committee, Spring Into
    Science science outreach, Princeton University (2024).
\end{itemize}

% ---------- Publications ----------
\section*{Selected First-Author Publications (7 first-author of 30+ total)}
\begin{itemize}
  \item Pouyaei, A.\ et al.\ (2026, under review). Linking Prognostic Fire
    Emissions to Atmospheric Chemistry and Aerosols in GFDL's Atmospheric and
    Land Model version 4.2 (AM4.2/LM4.2). \textit{J.\ Advances in Modeling
    Earth Systems.} \url{https://doi.org/10.22541/essoar.15005375/v1}
  \item Pouyaei, A.\ et al.\ (2025). ENSO-driven Variability in Ozone Sources
    and Its Impact on Tropospheric Ozone Radiative Forcing.
    \textit{Geophysical Research Letters}, 52(15).
  \item Pouyaei, A.\ et al.\ (2025). Implementation of Dynamic Fire Injection
    Height in GFDL's AM4.0: Impacts on Aerosol Profiles and Radiation.
    \textit{J.\ Advances in Modeling Earth Systems}, 17(4).
  \item Pouyaei, A.\ et al.\ (2023). Downwind Ozone Changes of the 2019
    Williams Flats Wildfire: WRF-Chem/DART Assimilation of OMI NO$_2$, HCHO,
    and MODIS AOD. \textit{J.\ Geophysical Research: Atmospheres}, 128(9).
  \item Pouyaei, A.\ et al.\ (2020). C-TRAIL Model v1.0 derived from CMAQ
    Model v5.2. \textit{Geoscientific Model Development}, 13(8).
\end{itemize}

\end{document}
